The guiding question of this work is the conservation of a generated
structure when its mathematical realization changes. Triadic Modular
Holography produces, through lifted arithmetic operations, states that
simultaneously retain orientation, incidence, boundary and memory. Its
discrete structure of the continuum assembles the compatible prolongations of
these states. The Archimedean realization yields coordinates through
cylinder refinement; the conjugate realization preserves exceptional
incidence relations. Genealogy thus precedes both the values
and the forms in which they are represented. This is the initial thesis that
the following formal development sets out with its operations and dependencies.

The article's specific geometric problem begins when a positive output
of that construction becomes a separation between nodes,
a network weight, a Plücker coordinate, a moment of a measure, an energy
coefficient or a direction of lattice deformation. Justifying this passage
requires determining what information each representation recovers, how
it compares with the others and what its refinement preserves. The significance
of a positive realization increases decisively when it admits a
marked inverse: geometry then ceases to provide only an evaluation
of the state and also preserves the information used by its readers.
The marking comprises the order of positions, the total scale and the
enriched data required by the particular reader.

The central contribution is organized around three operations. The first
is reconstruction: charge and memory recover the coordinates
of the record, and ratios of minors recover its separations.
The second is compatible refinement: a partition is subdivided
while preserving its endpoints, its markings and an exact record of the detail.
The third is transport of structures: a form, an operator or a
reader is transferred to another representation by an explicit map. These
operations establish conservation of residues, effective energy,
moments and spectra in the specific domains where each identity has
meaning. The adjective positive thus acquires several precise
mathematical realizations, articulated through reconstruction of the same datum.

Positive geometry has its own antecedents, which allow the level
of this contribution to be identified precisely. Boundary measurements
of directed networks parametrize cells of the totally nonnegative
Grassmannian in Postnikov's work~\cite{postnikov}; Scott describes the
cluster-algebra structure of the Grassmannian coordinate
ring~\cite{scott}. Arkani-Hamed and Trnka introduce the amplituhedron
in the geometric description of amplitudes in supersymmetric
\(\mathcal N=4\) Yang--Mills theory in the planar limit~\cite{amplituhedron}.
The formulation of positive geometries subsequently characterizes their canonical
forms by logarithmic poles and recursive boundary
residues~\cite{positive}. The present development composes these languages
of realization with data previously generated by HMT. Its distinctive
result is a marked, recoverable and refinable collection, together with
the maps linking its positive form and its material readings.

Refinement also extends the dimensional scope. The number of nodes
grows through nonadic subdivisions and permits construction of positive
matrices of every finite rank and external dimension compatible with that
number. The proof rests on Vandermonde determinants and full-rank
moment forms. In rank one, the complete convex geometry is obtained,
with simplicial coverage and residues in every finite dimension.
For higher ranks, an explicit collection of amplituhedral images
is constructed and its marked flags are preserved. Global coverage
of a higher-rank amplituhedron is a statement of a different scope
from the existence and compatibility theorem for this collection; the
subsequent proofs preserve that distinction in their formulation.

The physical dimension is introduced through the corresponding operators.
The tritic calendar distinguishes elliptic rotation, hyperbolic opening and
parabolic memory transport. Their Hamiltonian realizations preserve
the symplectic form and allow calculation of the Legendre transform and
the connection terms arising from a variable chart. In the gauge
sector, the areal--volumetric scale, common measure and reflection
forms precede reconstruction of the Yang--Mills Hamiltonian.
The spectral argument is presented together with uniform coercivity, the
witness attaining the first positive level and the intertwiners that
preserve it. The relation with canonical forms is established through their
readers and residues, while the gap belongs to the Hamiltonian and
its physical space. This hierarchy places each result where
its proof actually produces the conclusion.

The exposition first incorporates the formal construction required to
obtain the state and its record. It then develops positive
realizations, dimensional extension, energy and lattice duality;
finally, it brings together the Hamiltonian realizations and the gauge development.
Verification programs accompany specified identities and
explicit negative controls. Symbolic proof retains the
general scope; exact execution checks the declared specializations
and their correspondence with the printed formulas.
